\chapter{核心代码实现}\label{chap:code}

\section{BFS 核心代码}

BFS 的核心思想是使用队列按层遍历。代码如下：

\begin{lstlisting}[language=Python,caption={BFS 的核心实现},label={code:bfs}]
def bfs_search(grid, start, end):
    queue = deque([start])
    visited = {start}
    came_from = {}

    while queue:
        current = queue.popleft()
        if current == end:
            return reconstruct_path(came_from, current)

        for nxt in get_neighbors(grid, current):
            if nxt not in visited:
                visited.add(nxt)
                came_from[nxt] = current
                queue.append(nxt)
\end{lstlisting}

\section{A* 核心代码}

A* 通过优先队列维护待访问节点，并用 $f=g+h$ 来确定扩展优先级。

\begin{lstlisting}[language=Python,caption={A* 的节点扩展逻辑},label={code:astar}]
def a_star_search(grid, start, end):
    open_set = []
    heapq.heappush(open_set, (0, start))
    g_score = {start: 0}

    while open_set:
        _, current = heapq.heappop(open_set)
        if current == end:
            return True

        for nxt in get_neighbors(grid, current):
            tentative_g = g_score[current] + 1
            heuristic = abs(nxt[0] - end[0]) + abs(nxt[1] - end[1])
            f_score = tentative_g + heuristic
            heapq.heappush(open_set, (f_score, nxt))
\end{lstlisting}

\section{改进 A* 核心代码}

改进 A* 主要体现在代价函数设计上，代码如下：

\begin{lstlisting}[language=Python,caption={改进 A* 的附加代价},label={code:improved-astar}]
step_cost = 1.0
step_cost += nearby_obstacle_penalty(grid, nxt)
if prev_dir is not None and direction != prev_dir:
    step_cost += 0.1

heuristic = manhattan(nxt, end)
heuristic += 0.2 * nearby_obstacle_penalty(grid, nxt)
f_score = tentative_g + heuristic_weight * heuristic
\end{lstlisting}

\section{Q-Learning 核心代码}

Q-Learning 通过奖惩更新 Q 表，代码如下：

\begin{lstlisting}[language=Python,caption={Q-Learning 的价值更新},label={code:qlearning}]
reward = self.step_reward(state, nxt, valid, seen)
best_next = np.max(self.q_table[nxt[0], nxt[1]])
td_target = reward + self.gamma * best_next
td_error = td_target - self.q_table[state[0], state[1], action]
self.q_table[state[0], state[1], action] += self.alpha * td_error
\end{lstlisting}

\section{Dijkstra 核心代码}

Dijkstra 的重点是按当前最小路径代价不断扩展节点：

\begin{lstlisting}[language=Python,caption={Dijkstra 的核心实现},label={code:dijkstra}]
priority_queue = [(0.0, start)]
distances = {start: 0.0}

while priority_queue:
    current_cost, current = heapq.heappop(priority_queue)
    if current == end:
        return True
    for neighbor in get_neighbors(grid, current):
        new_cost = current_cost + 1.0
        if new_cost < distances.get(neighbor, float("inf")):
            distances[neighbor] = new_cost
            heapq.heappush(priority_queue, (new_cost, neighbor))
\end{lstlisting}

\section{动态可视化核心代码}

动态展示的关键是根据访问顺序和路径回溯过程逐帧更新图像：

\begin{lstlisting}[language=Python,basicstyle=\footnotesize\ttfamily,caption={动态展示中的帧更新逻辑},label={code:gif}]
def update(frame_idx):
    display = np.copy(base_grid)

    for node in visited_frames[:frame_idx]:
        if node not in (start, end) and display[node] == 0:
            display[node] = 2

    if frame_idx >= len(visited_frames) and result.success:
        path_idx = frame_idx - len(visited_frames)
        for node in path_frames[:path_idx]:
            if node not in (start, end):
                display[node] = 3

    image.set_data(display)
    return [image]
\end{lstlisting}
